\section{Conclusions}
\label{sec:concl}

\paragraph{Findings and contributions.}
A single hierarchical name, used unchanged by the registry, the metrics, the SIEM and the logs, was
the most useful design decision: every alert in this report names the pod, node and data centre
involved without any lookup. The same name made location-aware discovery and per-location
monitoring cheap to add, and federating only per-location aggregates reduced the central monitoring
traffic to 0.7\% of a raw federation. Wazuh detected all three threats and responded automatically
on both platforms, with no false positive against the legitimate clients. Over five repeated runs a
tampered DNS zone was restored after 0.28~s on average and a mail flood was blocked after 3.0~s;
flood sources were blocked within about 3~s in three runs, but in the other two the attacked
service was too overloaded to log the requests the detection depends on. Prometheus' correlation rules reduced each storage failure to one root-cause
alert. Moving the platform to Kubernetes required changes that only showed up under attack: a
block that acts before \code{kube-proxy}, an exception for the pod network, faster discovery of new
pod logs, and gentler liveness probes.

\paragraph{Why the solution is not perfect.}
\begin{itemize}
  \item \textbf{Detection depends on the victim.} Requests are logged when they complete, so a
        service too overloaded to answer is also too overloaded to be protected: in one of five
        DoS runs nothing was detected at all.
  \item \textbf{Per-source thresholds.} Both flood rules count requests per source address. In the
        test, 40 sources at one request per second each produced no Wazuh alert at all; only
        Prometheus saw the flood, and nothing responds automatically.
  \item \textbf{The SIEM can be blinded.} Each node's agent forwards at most about 500 events per
        second; above that, events from every pod on the node are lost or delayed by more than
        10~s. A flood therefore also hides whatever else happens on the attacked node.
  \item \textbf{Trusted interior.} Internal sources are exempt from the flood rule to avoid a
        self-inflicted outage, so a compromised pod can flood its neighbours unchallenged.
  \item \textbf{Unauthenticated registry.} A fake registration would redirect traffic like a DNS
        hijack and is invisible to file-integrity monitoring.
  \item \textbf{Narrow DNS coverage.} Only tampering with the zone file is detected; cache
        poisoning, rogue resolvers, on-path spoofing and CoreDNS changes are not.
  \item \textbf{Weak content filtering.} With untrained Bayes, SpamAssassin caught 70\% of a public
        spam corpus and would have quarantined 10\% of legitimate mail.
  \item \textbf{Hand-written correlation.} Root-cause rules exist per dependency (storage,
        ingestion); new services need new rules.
  \item \textbf{Single points of failure and a single host.} One registry, one SQLite writer, one
        Wazuh manager, privileged k3s nodes, and all locations on one machine: the numbers are lower
        bounds, the scalability measurements are noisy (Section~\ref{sec:repeat}), and the root of
        the deployment is a single failure domain.
  \item \textbf{Defences can hurt.} Twice the DoS response blocked the platform's own service, and the liveness
        probes once killed healthy pods under load. Every response needs its own false-positive
        test.
\end{itemize}

\paragraph{Future work.}
The second project (attack tolerance) can build directly on these gaps: authenticated, replicated
registries per location with the hierarchy delegated like DNS zones; replicated or sharded
storage; distributed-flood detection that feeds Prometheus anomalies into Wazuh so they can trigger
responses (for example rate limiting per location rather than per source); Kubernetes audit logs
and network policies to cover CoreDNS tampering and pod-to-pod abuse; a dependency graph or
Alertmanager inhibition to generate root-cause rules; a Wazuh manager cluster; and repeated
measurements on separate physical hosts.
